%% Part V: Why it scales. Two data structures that make waiting free, and the measurements to prove it.

\section{Inside epoll}
\label{sec:wakers-ch14}

\textit{A worker can sleep while ten thousand sockets wait, then wake for the one that matters, without examining the other 9,999. The reason lies in two lists the kernel keeps.}

\begin{description}
    \item[Interest list] Every registered descriptor, organized as a red-black tree. mio adds each socket once, edge-triggered, when it is created; Tokio never re-registers a socket for each read.
    \item[Ready list] A short linked list. When a packet lands in a socket's buffer, the kernel's networking code runs a small callback that appends just that socket, and wakes any thread sleeping in \code{epoll\_wait}.
\end{description}
\index{epoll}

\begin{figure}[H]
    \centering
    % Coordinates follow a 700 x 270 grid; y grows downwards.
    \begin{tikzpicture}[x=0.022cm, y=-0.022cm, font=\sffamily\scriptsize, >={Stealth}]
        % interest list
        \draw[rounded corners=6pt, draw=gray!50, fill=white, line width=0.6pt] (10,10) rectangle (280,260);
        \node[anchor=north west, inner sep=0] at (26,22) {\small\textbf{Interest list}};
        \node[anchor=north west, inner sep=0, text=gray] at (26,42) {every registered socket (red-black tree)};
        \foreach \r in {0,...,11}{
            \foreach \c in {0,...,13}{
                \fill[rounded corners=0.5pt, gray!25] ({26+17*\c},{68+14*\r}) rectangle ++(12,10);
            }
        }
        \foreach \c/\r in {3/2, 8/6, 12/9}{
            \fill[rounded corners=0.5pt, orange!70!black] ({26+17*\c},{68+14*\r}) rectangle ++(12,10);
        }
        \node[anchor=west, inner sep=0, text=gray] at (26,246) {registered once, at creation $\cdot$ mostly idle};
        % packet arrives
        \draw[->, gray, thick] (282,122) -- (386,122);
        \node[text=orange!70!black, align=center] at (334,102) {packet arrives:};
        \node[text=orange!70!black, align=center] at (334,142) {callback appends};
        % ready list
        \draw[rounded corners=6pt, draw=orange!70!black, fill=yellow!10, line width=0.6pt] (390,62) rectangle (524,182);
        \node[anchor=north west, inner sep=0] at (404,74) {\small\textbf{Ready list}};
        \foreach \fd [count=\k from 0] in {912, 4817, 7310}{
            \fill[rounded corners=2pt, white] (404,{98+26*\k}) rectangle ++(106,20);
            \node[anchor=west, font=\ttfamily\scriptsize, text=orange!70!black] at (412,{108+26*\k}) {fd \fd};
        }
        \draw[->, gray, thick] (524,122) -- (564,122);
        % epoll_wait
        \draw[rounded corners=6pt, draw=StmtLemma, fill=green!6, line width=0.6pt] (568,62) rectangle (690,182);
        \node[anchor=north west, inner sep=0, text width=2.3cm] at (580,74) {\small\textbf{epoll\_wait}\\[3pt] \scriptsize\textcolor{StmtLemma}{returns 3 events}\\[6pt] driver calls\\ 3 $\times$ \texttt{wake()}};
        \node[anchor=north west, inner sep=0, text=gray, text width=6.6cm] at (390,204) {Work done is proportional to the 3 ready sockets, not to the thousands being watched.};
    \end{tikzpicture}
    \caption{epoll's two lists. Highlighted squares are the sockets that just received data.}
    \label{fig:wakers-14.1}
\end{figure}

\begin{lemma}{Event-proportional waiting}{event-proportional-waiting}
    The work done by one call to \code{epoll\_wait} is proportional to the number of ready descriptors it returns, not to the number registered.
    \why Descriptors join the ready list as their events happen, through the kernel's callback. \code{epoll\_wait} only copies out the ready list; it never walks the interest list.
\end{lemma}

\begin{table}[H]
    \centering\small
    \begin{tabular}{r c r}
        \toprule
        \textbf{Sockets watched} & \textbf{Entries \texttt{epoll\_wait} touches} & \textbf{Entries \texttt{poll()}/\texttt{select()} scan} \\
        \midrule
        1,000 & \textcolor{StmtLemma}{\textbf{3}} & \textcolor{red!70!black}{\textbf{1,000}} \\
        10,000 & \textcolor{StmtLemma}{\textbf{3}} & \textcolor{red!70!black}{\textbf{10,000}} \\
        100,000 & \textcolor{StmtLemma}{\textbf{3}} & \textcolor{red!70!black}{\textbf{100,000}} \\
        1,000,000 & \textcolor{StmtLemma}{\textbf{3}} & \textcolor{red!70!black}{\textbf{1,000,000}} \\
        \bottomrule
    \end{tabular}
    \caption{One round in which three sockets receive data.}
    \label{tab:wakers-14.1}
\end{table}

\begin{quotebox}{Summary}
    \begin{itemize}
        \item The interest list holds everything; the ready list holds only what changed.
        \item Waiting costs work in proportion to events, not to watchers (Lemma~\ref{lem:event-proportional-waiting}). A million idle sockets cost kernel memory, but no CPU.
    \end{itemize}
\end{quotebox}

\section{Inside the timer wheel}
\label{sec:wakers-ch15}

\textit{Timers in Tokio are not a kernel feature. Every \code{sleep} and \code{timeout} lives in a single user-space structure built like a clock with six hands, and its next deadline becomes \code{epoll\_wait}'s timeout.}

\begin{definition}{Hierarchical timer wheel}{timer-wheel}
    \index{Timer wheel}
    Tokio's timer wheel has six levels of 64 slots. A level-0 slot spans 1 millisecond, and each level above is 64 times coarser: 64\,ms, about 4.1\,s, 4.4\,min, 4.7\,h and 12.4 days. Together the levels span $64^6$ milliseconds, a little over two years. Each level also keeps a 64-bit mask recording which of its slots are occupied.
\end{definition}

\begin{lemma}{Constant-time insert}{constant-time-insert}
    Adding a timer takes constant time, however many timers exist.
    \why The level is chosen from the highest 6-bit group in which ``now'' and the deadline differ; the slot from the deadline's bits at that level. The timer is appended to that slot's list. None of this depends on the number of timers.
\end{lemma}

\begin{lemma}{Constant-time next deadline}{constant-time-next-deadline}
    Finding the next deadline takes at most one bit operation per level: six in all.
    \why On each level, a rotate of the occupied-slot mask followed by a count of trailing zeros finds the next non-empty slot.
\end{lemma}

\begin{lemma}{Timers are timeouts}{timers-are-timeouts}
    A timer fires when the driver runs at or after its deadline and calls the wakers in the due slot. On an idle runtime this happens because \code{epoll\_wait}'s timeout expires; on a busy one, a worker's periodic non-blocking check of the driver (Table~\ref{tab:wakers-11.1}) can also do it.
    \why The time driver passes ``time until the next deadline'' (Lemma~\ref{lem:constant-time-next-deadline}) as the timeout (Listing~\ref{lst:wakers-11.1}), and after every return from the layers below it, processes whatever timers are due. Timers have no kernel object of their own.
\end{lemma}

Coarse timers \emph{cascade}: when the window of a higher-level slot arrives, its timers are re-filed into finer levels, and they finally fire from level 0.

\subsection{A worked example}

At time zero, five timers are created, with deadlines of 3\,ms, 70\,ms, 300\,ms, 5\,s and 90\,s. Figure~\ref{fig:wakers-15.1} shows where Tokio's placement rule puts them.

\begin{figure}[H]
    \centering
    % Coordinates follow a 774 x 308 grid; y grows downwards. Slot i of a level starts at x = 118 + 10i.
    \newcommand{\wklevel}[4]{% top y, name, slot size, occupied slots
        \node[anchor=north west, inner sep=0] at (8,{#1+6}) {\small\textbf{#2}};
        \node[anchor=north west, inner sep=0, text=gray] at (8,{#1+22}) {#3};
        \foreach \t in {0,16,32,48,63}{\node[text=gray, font=\sffamily\tiny] at ({122+10*\t},{#1-6}) {\t};}
        \foreach \s in {0,...,63}{\fill[rounded corners=0.5pt, gray!15] ({118+10*\s},#1) rectangle ++(8,26);}
        \foreach \s in {#4}{\fill[rounded corners=0.5pt, StmtLemma] ({118+10*\s},#1) rectangle ++(8,26);}
        \draw[StmtProposition, line width=0.9pt, rounded corners=1pt] (116,{#1-2}) rectangle ++(12,30);
    }
    \begin{tikzpicture}[x=0.02cm, y=-0.02cm, font=\sffamily\scriptsize]
        \wklevel{20}{Level 0}{64 $\times$ 1 ms}{3}
        \wklevel{116}{Level 1}{64 $\times$ 64 ms}{1,4}
        \wklevel{212}{Level 2}{64 $\times$ 4.1 s}{1,21}
        % timer labels
        \draw[StmtLemma] (152,47) -- (152,58);
        \node[anchor=north, align=center, inner sep=1pt] at (152,58) {\textcolor{StmtLemma}{\textbf{3 ms}}\\ \textcolor{gray}{\tiny slot 3}};
        \draw[StmtLemma] (132,143) -- (132,154);
        \node[anchor=north, align=center, inner sep=1pt] at (132,154) {\textcolor{StmtLemma}{\textbf{70 ms}}\\ \textcolor{gray}{\tiny slot 1}};
        \draw[StmtLemma] (162,143) -- (162,150) -- (180,150);
        \node[anchor=west, inner sep=1pt] at (180,150) {\textcolor{StmtLemma}{\textbf{300 ms}} \textcolor{gray}{\tiny $\cdot$ slot 4}};
        \draw[StmtLemma] (132,239) -- (132,250);
        \node[anchor=north, align=center, inner sep=1pt] at (132,250) {\textcolor{StmtLemma}{\textbf{5 s}}\\ \textcolor{gray}{\tiny slot 1}};
        \draw[StmtLemma] (332,239) -- (332,250);
        \node[anchor=north, align=center, inner sep=1pt] at (332,250) {\textcolor{StmtLemma}{\textbf{90 s}}\\ \textcolor{gray}{\tiny slot 21}};
    \end{tikzpicture}
    \caption{Five timers placed at time zero. The first three of the six levels are shown; the orange outline marks the current slot on each level.}
    \label{fig:wakers-15.1}
\end{figure}

Table~\ref{tab:wakers-15.1} follows time forward. Each row is one return from \code{epoll\_wait}, and each touches only the single slot that came due.

\begin{table}[H]
    \centering\small
    \begin{tabularx}{\linewidth}{r r X}
        \toprule
        \textbf{Time} & \textbf{Timeout} & \textbf{What the driver does} \\
        \midrule
        3 ms & 3 ms & Level-0 slot 3 due: \textbf{fire} the 3 ms timer \\
        64 ms & 61 ms & Level-1 slot 1 due: cascade the 70 ms timer into level-0 slot 6 \\
        70 ms & 6 ms & \textbf{Fire} the 70 ms timer \\
        256 ms & 186 ms & Level-1 slot 4 due: cascade the 300 ms timer into level-0 slot 44 \\
        300 ms & 44 ms & \textbf{Fire} the 300 ms timer \\
        4,096 ms & 3,796 ms & Level-2 slot 1 due: cascade the 5 s timer into level-1 slot 14 \\
        4,992 ms & 896 ms & Level-1 slot 14 due: cascade it into level-0 slot 8 \\
        5,000 ms & 8 ms & \textbf{Fire} the 5 s timer \\
        86,016 ms & 81,016 ms & Level-2 slot 21 due: cascade the 90 s timer into level-1 slot 62 \\
        89,984 ms & 3,968 ms & Cascade it into level-0 slot 16 \\
        90,000 ms & 16 ms & \textbf{Fire} the 90 s timer \\
        \bottomrule
    \end{tabularx}
    \caption{The example run to completion, computed with Tokio's rule. A run of 10,030 random timers, checked the same way, fired every one exactly on its deadline.}
    \label{tab:wakers-15.1}
\end{table}

\begin{itemize}
    \item Ninety seconds of timekeeping for five timers costs eleven brief wake-ups.
    \item Between them, the worker sleeps (Axiom~\ref{ax:free-blocking}).
    \item Adding ten thousand more timers adds work only at the moments when some of them are due.
\end{itemize}

\begin{quotebox}{Summary}
    \begin{itemize}
        \item Six levels of 64 slots (Definition~\ref{def:timer-wheel}); constant-time insert and next-deadline (Lemmas~\ref{lem:constant-time-insert}, \ref{lem:constant-time-next-deadline}).
        \item Timers cascade from coarse to fine levels, and fire as \code{epoll\_wait} timeouts (Lemma~\ref{lem:timers-are-timeouts}).
    \end{itemize}
\end{quotebox}

\section{Idle cost and event cost}
\label{sec:wakers-ch16}

\textit{A waker is not a thread, and a waiting task runs nothing. That claim follows from the axioms; this section states it precisely and then measures it.}

\begin{definition}{Idle cost, event cost}{idle-event-cost}
    \index{Idle cost}\index{Event cost}
    The \emph{idle cost} of a waiting task is what it consumes while nothing happens to it. The \emph{event cost} is the work done when its wake source fires.
\end{definition}

\begin{lemma}{Idle cost is memory only}{idle-cost-memory-only}
    On a waker-driven runtime, the idle cost of a task is memory---the task's allocation and its waker slot---and no CPU time and no thread.
    \why A waiting task's waker sits in a waker slot (Definition~\ref{def:waker-slot}), which is data. The task is not polled (Definition~\ref{def:waker-driven-executor}). Its wake source is watched by the kernel or a timer wheel, and the thread that waits on them is blocked, which costs no CPU (Axiom~\ref{ax:free-blocking}).
\end{lemma}

\begin{lemma}{Event cost is independent of waiters}{event-cost-independent}
    The work done when one task's wake source fires does not grow with the number of other waiting tasks.
    \why epoll returns only ready descriptors (Lemma~\ref{lem:event-proportional-waiting}); the timer wheel touches only the due slot (Lemmas~\ref{lem:constant-time-insert}, \ref{lem:constant-time-next-deadline}); the token leads straight to the waker slot (Lemma~\ref{lem:token-round-trip}); and the run queue holds only woken tasks.
\end{lemma}

\subsection{Measurement}

The test ran on a two-core Linux machine with Tokio 1.54.1, a four-worker runtime and a release build (Listing~\ref{lst:wakers-A.4}):

\begin{enumerate}
    \item spawn 100,000 tasks that each sleep for over ten minutes;
    \item open 9,000 loopback TCP connections, each with a task waiting in \code{read().await};
    \item after each stage, sample the process during five idle seconds;
    \item finally, write one byte to one more socket and time its task's wake-up.
\end{enumerate}

\begin{table}[H]
    \centering\small
    \begin{tabular}{l r r r}
        \toprule
        \textbf{State} & \textbf{Resident memory} & \textbf{OS threads} & \textbf{CPU used in 5 s} \\
        \midrule
        Empty runtime & 2.9 MiB & 5 & 0 ms \\
        + 100,000 sleeping tasks & 39.7 MiB & 5 & 0 ms \\
        + 9,000 idle sockets (18,000 descriptors) & 49.9 MiB & 5 & 0 ms \\
        \bottomrule
    \end{tabular}
    \caption{Idle cost. ``0 ms'' means below the 10 ms resolution of the kernel's CPU accounting. The kernel's own socket memory is not part of the resident figure.}
    \label{tab:wakers-16.1}
\end{table}

The results match the lemmas:

\begin{itemize}
    \item \textbf{Lemma~\ref{lem:idle-cost-memory-only}:} about 380 bytes per sleeping task including its timer entry, and about 1.1\,KB per socket task in user space. The thread count never changed, and idle CPU was unmeasurable.
    \item \textbf{Lemma~\ref{lem:event-cost-independent}:} the final write woke its task in 84 microseconds, with 109,000 other waiters in the same runtime.
\end{itemize}

The real limits on an async server therefore lie elsewhere:

\begin{itemize}
    \item memory;
    \item the per-process file-descriptor limit (\code{ulimit -n});
    \item how much work each event triggers;
    \item tasks that compute for a long time without reaching an \code{.await}, since tasks share threads cooperatively and such a task stalls every other task on its worker.
\end{itemize}

\begin{quotebox}{Summary}
    \begin{itemize}
        \item Idle cost is memory only (Lemma~\ref{lem:idle-cost-memory-only}). Event cost does not depend on the number of waiters (Lemma~\ref{lem:event-cost-independent}).
        \item Measured: 109,000 waiters in about 50 MiB, with no extra threads, no idle CPU, and an 84\,$\mu$s wake-up.
    \end{itemize}
\end{quotebox}
